% GENERATED FILE. Edit canonical lessons, then run npm run generate:books.
% canonical-source-hash: fnv1a64:8053ff400f978e1c
% canonical-lessons: MR-C308-pidhi, MR-C308-purvaj, MR-C308-varsa, MR-C308-vams, MR-C308-jatan-karne

\chapter{Generations}
\label{ch:mr-generations}
\hlchaptermodality{308}

\begin{chapteropening}
\textbf{By the end of this chapter:} \emph{I can talk about ancestors, heritage and what a family keeps.}

Complete the last lesson of chapter 308: I can talk about ancestors, heritage and what a family keeps.
\end{chapteropening}

\section[\texorpdfstring{\mr{पिढी} (piḍhī)}{पिढी (piḍhī)}]{\mr{पिढी} (piḍhī) — a generation}

\label{lesson:MR-C308-pidhi}



\begin{quote}
Before the new one: say the Marathi for a period, then the Marathi for a remembrance.
\end{quote}

\subsection*{You'll want to know: \mr{पिढी}}
\textbf{\mr{पिढी}} — \emph{piḍhī} — \textquotedblleft{}a generation\textquotedblright{}.

\begin{quote}
\mr{आमच्या} \mr{घरात} \mr{तीन} \mr{पिढ्या} \mr{एकत्र} \mr{राहतात}.
\end{quote}

\begin{quote}
\emph{āmacyā gharāt tīn piḍhyā ekatra rāhtāt.}
\end{quote}

\begin{quote}
Three generations live together in our house.
\end{quote}

Feminine. Plural: \mr{पिढ्या}, \emph{piḍhyā}.

Plural \mr{पिढ्या}; \mr{पुढची} \mr{पिढी}, the next generation.

\begin{cousinweb}
\textbf{\mr{पिढी}} (\emph{piḍhī}) is the Marathi word for a generation; its origin is not settled.
\end{cousinweb}

\begin{grammarlens}[title={What you've built}]
One more word for telling what happened.
\end{grammarlens}

\subsection*{Your turn}
\begin{itemize}
\raggedright
  \item \emph{Say it:} \emph{piḍhī}
  \item \emph{Say it:} \emph{piḍhī}, once more
  \item \emph{Say it:} \emph{āmacyā gharāt tīn piḍhyā ekatra rāhtāt.} — Three generations live together in our house.
  \item \emph{Recall:} say the Marathi for a period, then the Marathi for a remembrance, then say \emph{piḍhī} again
\end{itemize}

\begin{tcolorbox}[breakable,colback=teal!4,colframe=teal!35!black,title={Before you move on}]
What does \mr{पिढी} mean? (\textquotedblleft{}A generation\textquotedblright{}.) Say it once more.
\end{tcolorbox}
\section[\texorpdfstring{\mr{पूर्वज} (pūrvaj)}{पूर्वज (pūrvaj)}]{\mr{पूर्वज} (pūrvaj) — an ancestor}

\label{lesson:MR-C308-purvaj}



\begin{quote}
Before the new one: say the Marathi for a remembrance, then the Marathi for a generation.
\end{quote}

\subsection*{You'll want to know: \mr{पूर्वज}}
\textbf{\mr{पूर्वज}} — \emph{pūrvaj} — \textquotedblleft{}an ancestor\textquotedblright{}.

\begin{quote}
\mr{आमचे} \mr{पूर्वज} \mr{या} \mr{गावात} \mr{तीनशे} \mr{वर्षांपूर्वी} \mr{आले}.
\end{quote}

\begin{quote}
\emph{āmce pūrvaj yā gāvāt tīnśe varṣāmpūrvī āle.}
\end{quote}

\begin{quote}
Our ancestors came to this village three hundred years ago.
\end{quote}

Masculine.

Mostly in the plural: \mr{पूर्वजांचं} \mr{घर}, the ancestral house.

\begin{cousinweb}
\textbf{\mr{पूर्वज}} (\emph{pūrvaj}) is a tatsama, taken from Sanskrit \emph{pūrvaja}, 'born before; an ancestor' (Monier-Williams).
\end{cousinweb}

\begin{grammarlens}[title={What you've built}]
One more word for telling what happened.
\end{grammarlens}

\subsection*{Your turn}
\begin{itemize}
\raggedright
  \item \emph{Say it:} \emph{pūrvaj}
  \item \emph{Say it:} \emph{pūrvaj}, once more
  \item \emph{Say it:} \emph{āmce pūrvaj yā gāvāt tīnśe varṣāmpūrvī āle.} — Our ancestors came to this village three hundred years ago.
  \item \emph{Recall:} say the Marathi for a remembrance, then the Marathi for a generation, then say \emph{pūrvaj} again
\end{itemize}

\begin{tcolorbox}[breakable,colback=teal!4,colframe=teal!35!black,title={Before you move on}]
What does \mr{पूर्वज} mean? (\textquotedblleft{}An ancestor\textquotedblright{}.) Say it once more.
\end{tcolorbox}
\section[\texorpdfstring{\mr{वारसा} (vārsā)}{वारसा (vārsā)}]{\mr{वारसा} (vārsā) — heritage; an inheritance}

\label{lesson:MR-C308-varsa}



\begin{quote}
Before the new one: say the Marathi for a generation, then the Marathi for an ancestor.
\end{quote}

\subsection*{You'll want to know: \mr{वारसा}}
\textbf{\mr{वारसा}} — \emph{vārsā} — \textquotedblleft{}heritage; an inheritance\textquotedblright{}.

\begin{quote}
\mr{हा} \mr{किल्ला} \mr{आपला} \mr{ऐतिहासिक} \mr{वारसा} \mr{आहे}.
\end{quote}

\begin{quote}
\emph{hā killā āplā aitihāsik vārsā āhe.}
\end{quote}

\begin{quote}
This fort is our historical heritage.
\end{quote}

Masculine.

\mr{वारसा} \mr{जपणे}, to cherish a heritage.

\begin{cousinweb}
\textbf{\mr{वारसा}} (\emph{vārsā}) is built on \emph{vāras}, 'an heir', a loan from Persian \emph{vāris}, usually traced to Arabic \emph{wārith}, 'an heir'.
\end{cousinweb}

\begin{grammarlens}[title={What you've built}]
One more word for telling what happened.
\end{grammarlens}

\subsection*{Your turn}
\begin{itemize}
\raggedright
  \item \emph{Say it:} \emph{vārsā}
  \item \emph{Say it:} \emph{vārsā}, once more
  \item \emph{Say it:} \emph{hā killā āplā aitihāsik vārsā āhe.} — This fort is our historical heritage.
  \item \emph{Recall:} say the Marathi for a generation, then the Marathi for an ancestor, then say \emph{vārsā} again
\end{itemize}

\begin{tcolorbox}[breakable,colback=teal!4,colframe=teal!35!black,title={Before you move on}]
What does \mr{वारसा} mean? (\textquotedblleft{}Heritage; an inheritance\textquotedblright{}.) Say it once more.
\end{tcolorbox}
\section[\texorpdfstring{\mr{वंश} (vaṁś)}{वंश (vaṁś)}]{\mr{वंश} (vaṁś) — a lineage; a dynasty}

\label{lesson:MR-C308-vams}



\begin{quote}
Before the new one: say the Marathi for an ancestor, then the Marathi for heritage.
\end{quote}

\subsection*{You'll want to know: \mr{वंश}}
\textbf{\mr{वंश}} — \emph{vaṁś} — \textquotedblleft{}a lineage; a dynasty\textquotedblright{}.

\begin{quote}
\mr{त्या} \mr{घराण्याचा} \mr{वंश} \mr{अजून} \mr{चालू} \mr{आहे}.
\end{quote}

\begin{quote}
\emph{tyā gharāṇyācā vaṁś ajūn cālū āhe.}
\end{quote}

\begin{quote}
That family's line still continues.
\end{quote}

Masculine.

\mr{राजवंश}, a royal dynasty.

\begin{cousinweb}
\textbf{\mr{वंश}} (\emph{vaṁś}) is a tatsama, taken from Sanskrit \emph{vaṁśa}, 'a bamboo; a lineage' (Monier-Williams): a family line grows joint by joint, like bamboo.
\end{cousinweb}

\begin{grammarlens}[title={What you've built}]
One more word for telling what happened.
\end{grammarlens}

\subsection*{Your turn}
\begin{itemize}
\raggedright
  \item \emph{Say it:} \emph{vaṁś}
  \item \emph{Say it:} \emph{vaṁś}, once more
  \item \emph{Say it:} \emph{tyā gharāṇyācā vaṁś ajūn cālū āhe.} — That family's line still continues.
  \item \emph{Recall:} say the Marathi for an ancestor, then the Marathi for heritage, then say \emph{vaṁś} again
\end{itemize}

\begin{tcolorbox}[breakable,colback=teal!4,colframe=teal!35!black,title={Before you move on}]
What does \mr{वंश} mean? (\textquotedblleft{}A lineage; a dynasty\textquotedblright{}.) Say it once more.
\end{tcolorbox}
\section[\texorpdfstring{\mr{जतन} \mr{करणे} (jatan karṇe)}{जतन करणे (jatan karṇe)}]{\mr{जतन} \mr{करणे} (jatan karṇe) — to preserve}

\label{lesson:MR-C308-jatan-karne}



\begin{quote}
Before the new one: say the Marathi for heritage, then the Marathi for a lineage.
\end{quote}

\subsection*{You'll want to know: \mr{जतन} \mr{करणे}}
\textbf{\mr{जतन} \mr{करणे}} — \emph{jatan karṇe} — \textquotedblleft{}to preserve\textquotedblright{}.

\begin{quote}
\mr{संग्रहालयात} \mr{जुनी} \mr{नाणी} \mr{जतन} \mr{केली} \mr{आहेत}.
\end{quote}

\begin{quote}
\emph{saṅgrahālayāt junī nāṇī jatan kelī āhet.}
\end{quote}

\begin{quote}
Old coins have been preserved in the museum.
\end{quote}

Past, with -\mr{ने} (he …): \mr{त्याने} \mr{जतन} \mr{केलं}, \emph{tyāne jatan kelaṁ}.

Formal; what is kept agrees: \mr{परंपरा} \mr{जतन} \mr{केली}.

\begin{cousinweb}
\textbf{\mr{जतन} \mr{करणे}} (\emph{jatan karṇe}) joins \emph{jatan}, 'careful keeping', to \emph{karṇe}, 'to do'. \emph{jatan} is usually traced to Sanskrit \emph{yatna}, 'effort' (Monier-Williams), which Prakrit has as \emph{jatta} (CDIAL 10404a); CDIAL gives no Marathi form.
\end{cousinweb}

\begin{grammarlens}[title={What you've built}]
One more word for telling what happened.
\end{grammarlens}

\subsection*{Your turn}
\begin{itemize}
\raggedright
  \item \emph{Say it:} \emph{jatan karṇe}
  \item \emph{Say it:} \emph{jatan karṇe}, once more
  \item \emph{Say it:} \emph{saṅgrahālayāt junī nāṇī jatan kelī āhet.} — Old coins have been preserved in the museum.
  \item \emph{Recall:} say the Marathi for heritage, then the Marathi for a lineage, then say \emph{jatan karṇe} again
\end{itemize}

\begin{tcolorbox}[breakable,colback=teal!4,colframe=teal!35!black,title={Before you move on}]
What does \mr{जतन} \mr{करणे} mean? (\textquotedblleft{}To preserve\textquotedblright{}.) Say it once more.
\end{tcolorbox}
